\chapter{CONCLUSION}

This report presented a \textbf{Propositional Logic Evaluator} that applies classical compiler front-end techniques to evaluate expressions over truth values (\texttt{t}, \texttt{f}) with conjunction ($\wedge$) and disjunction ($\vee$). The system follows a four-stage pipeline --- lexical analysis, LALR(1) parsing (via SLY~\citep{beazley2023sly}), post-order AST evaluation, and pre-order prefix translation --- producing both the evaluated truth value and an equivalent prefix notation string from a single parse~\citep{aho2006compilers}. An interactive PySide6 GUI with dual Unicode/ASCII input modes completes the implementation. All six required test cases pass correctly (Chapter~4, Table~\ref{table:test_results}).

Three design choices merit emphasis: using an \textbf{LALR(1) parser generator} keeps the grammar declarative and aligned with formal BNF notation; constructing an \textbf{AST as an intermediate representation} decouples parsing from downstream processing, enabling independent tree traversals for evaluation and translation without re-parsing; and producing \textbf{fully parenthesised prefix notation} ensures unambiguous output regardless of operator precedence. Future extensions could include additional connectives ($\neg$, $\rightarrow$, $\leftrightarrow$), named propositional variables, automatic truth table generation, and step-by-step evaluation visualisation.
